% Generated by `taskgen course examples` from dlp.collision_probability (seed 4).
% Do not edit: change examples.yaml and regenerate.
\begin{example}[длина списков для DLP]\label{ex:collision-lists}
Решаем DLP в группе простого порядка $N$ поиском коллизии списков $g^{y_i}$ и $h \cdot g^{z_j}$ со случайными показателями. Какой длины $n$ взять два списка (одинаковой длины) случайных элементов множества из $N = 912\,049\,679$ элементов, чтобы по теореме о коллизиях вероятность общего элемента была не меньше $0{,}5$? Вычисления провести с логарифмами (4 знака после запятой).

\par\noindent\textit{Решение.} При $n = m$: $1 - e^{-n^2/N} \ge P \iff n^2 \ge N \ln\frac{1}{1 - P} \iff n \ge \sqrt{N \ln\frac{1}{1-P}}$. $\ln\frac{1}{1 - 0{,}5} = \ln 2 \approx 0{,}6931$, $N \ln\frac{1}{1-P} \approx 632\,184\,663{,}53$, $\sqrt{632\,184\,663{,}53} \approx 25\,143{,}28$. Проверка: $1 - e^{-25\,144^2/N} \approx 0{,}5000 \ge 0{,}5$, а при $n = 25\,143$: $\approx 0{,}499993 < 0{,}5$ (для сравнения, $\sqrt N \approx 30\,200{,}16$).\[n = m = 25\,144.\]
\end{example}
